%************************************************
\chapter{Introduction}\label{ch:introduction}
%************************************************

In recent years, generative artificial intelligence has fundamentally transformed the landscape of synthetic data generation. 
Early breakthroughs with Generative Adversarial Networks (GANs) and Variational Autoencoders (VAEs) demonstrated the potential 
of neural networks to learn and sample from complex data distributions. More recently, the advent of score-based Diffusion Models 
and Flow Matching frameworks has set a new standard, achieving unprecedented fidelity and training stability when synthesizing 
highly realistic, high-dimensional data. \\
However, as these powerful models are increasingly deployed in high-stakes domains such as healthcare, their inherently opaque, 
"black box" nature poses significant challenges. This thesis, titled Beyond the Black Box: Disentangled and Controllable Generative 
Models for Medical Imaging, aims to bridge the gap between raw generative power and clinical reliability. The primary focus of 
this research is to enhance the controllability, disentanglement, and compositionality of state-of-the-art generative frameworks—specifically 
Flow Matching and Diffusion models—for the generation of safe and reliable synthetic medical data. By establishing a deep theoretical 
understanding of these models' underlying principles, this work identifies their current structural limitations and proposes innovative, 
modular solutions to holistically address the complexities and severe scarcity of annotated medical datasets. \\
\\
The development and validation of robust medical AI systems require massive amounts of diverse, high-quality, and richly annotated data.
However, the acquisition of such datasets is chronically hindered by strict privacy regulations, high annotation costs, and the inherent 
rarity of certain pathologies. While synthetic data generation offers a promising solution to this bottleneck, current standard generative 
models often lack the precise control necessary for clinical applications. \\
When a generative model operates as a "black box," users cannot 
reliably predict how specific anatomical structures or pathological features will be rendered. The widespread deployment of medical AI hinges 
fundamentally on the concept of human-AI trust. However, as recent frameworks on AI trustworthiness emphasize, simply providing mathematical 
interpretability—trying to understand how the model works internally—is neither necessary nor sufficient for establishing this trust \cite{shen2022trustaiinterpretabilitynecessary}. 
Instead, human trust in AI is built through interaction. Users must be able to run the model at will and probe its boundaries to observe 
predictable outcomes, generating empirical "behavior certificates" that prove the model's reliability under specific constraints. \\
Therefore, to be truly safe and useful in a medical context, synthetic generation must prioritize this interactive capability through disentanglement. 
It must allow a clinician or researcher to alter a single semantic feature—such as the size, severity, or location of a tumor—and immediately observe 
the model's response without the surrounding healthy anatomy becoming distorted. The motivation for this research stems from the critical need to 
facilitate this interactive trust. By developing frameworks where different semantic components can be independently manipulated and recombined, 
this thesis strives to transition generative models from static "black boxes" into highly interactive, trustworthy tools that users can rigorously 
test, understand through behavior, and safely use to simulate rare clinical scenarios. \\
\\
To achieve the overarching goal of developing highly controllable and disentangled generative models for medical imaging, this research is driven 
by the following specific objectives:
\begin{itemize}
    \item \textbf{Theoretical Grounding}: Acquire an in-depth theoretical understanding of diffusion and flow-based generative models, with a specific focus 
    on their probabilistic foundations, interpolation strategies, and current structural limitations. 
    \item \textbf{Comprehensive Literature Review}: Conduct a thorough analysis of existing literature surrounding control mechanisms and compositional frameworks in 
    generative modeling, highlighting gaps in their application to the medical domain.
    \item \textbf{Baseline Benchmarking}: Reproduce and systematically benchmark baseline models, including Denoising Diffusion Probabilistic Models (DDPMs) and Flow 
    Matching models, utilizing public and synthetic medical datasets to quantitatively evaluate current limitations in controllability and fidelity. 
    \item \textbf{Design of Novel Control Strategies}: Propose and implement novel architectural modifications and training mechanisms—such as semantic conditioning, 
    class guidance, and spatial control—to enforce strict user-defined constraints on the generative process.
    \item \textbf{Compositional Generation}: Develop compositional generation techniques that enable modular synthesis, ensuring that distinct semantic components 
    (e.g., pathology presence and anatomical structure) can be independently disentangled, manipulated, and seamlessly recombined.
    \item \textbf{Clinical Validation and Evaluation}: Apply the proposed models to real-world medical datasets, rigorously evaluating the methods through both quantitative 
    metrics (e.g., FID, precision-recall, segmentation overlap) and qualitative assessments by clinical experts to ensure the utility, diversity, and 
    realism of the generated data. 
    \item \textbf{Dissemination and Open Science}: Contribute to the broader research community by disseminating findings through peer-reviewed publications 
    and providing open-source implementations of the developed frameworks to ensure transparency and reproducibility.
\end{itemize}

%*****************************************
%*****************************************
%*****************************************
%*****************************************
%*****************************************




